\section{Introduction}\label{sec:intro}

Submissions to arXiv more than doubled between 2018 and 2025, and in September
2026 the archive received a record 40,363 papers. Large language models now
draft text faster than any researcher can, and they are starting to read and
summarise for us too. Two opposite hopes and fears follow. Faster writing
might flood science with duplicates. Better reading tools might be the remedy
and raise research productivity. This paper argues that free choice of
research topics turns both intuitions around. Faster writing scales science,
while better reading tools are largely competed away.

The argument rests on one observation: researchers choose not only how much to
read and write but \emph{where} to work. Reading has two jobs. It
\emph{teaches}: what one reads raises the quality of one's own work, and more so
in a topic with more to learn from. It \emph{protects}: each earlier paper may
already cover one's project, and reading is how one finds out. Unaware
duplicates lose part of their credit; detected overlaps can be turned into
follow-ons that cite the earlier work, or abandoned. These two jobs make the
payoff of working in a topic rise with the flow of its literature while there is
much to learn, and fall once more of that flow goes unread. Newcomers enter
topics until they are no better than the frontier, where new topics are opened.

\paragraph{The headline result.} A tool that speeds up writing raises the payoff
of everything a researcher does, including opening a new topic. Newcomers spread
out, so output per researcher rises almost one for one with writing speed, and
the duplication rate barely moves (\cref{prop:ai,prop:solow}). A tool that
speeds up reading, search or catching up on a literature raises the payoff only
of working inside existing topics. Newcomers crowd in until those topics are
again no better than the frontier. Density rises, fewer topics are opened, and
output per researcher rises only to the extent that the frontier improves as
fewer directions are taken. The gain from better reading is dissipated by entry
in the same way that a better road into a city is capitalised into its size.
This asymmetry gives one reason why decades of better access to the literature
have coincided with falling measured research productivity
\citep{bloometal2020}.

\paragraph{What is new, and what is not.} The equilibrium structure is that of
a system of cities \citep{henderson1974} or of clubs \citep{buchanan1965}. The
average payoff is single peaked in group size, new groups enter freely,
equilibria are stable only on the falling side, the planner sits at the peak
of the average payoff, and an entry tax decentralises the optimum. Self-organised
cities are too large for the same reason \citep{hendersonbecker2000,scotchmer2002}.
We use this structure and claim four things that it does not provide. First,
a microfoundation of the hump in the two jobs of reading, valid for any
learning and detection technology in a broad class. The proof uses joint
log-concavity of the researcher's problem in reading time and density
(\cref{thm:v}), and the hump comes from bounded absorption against a covering
hazard that grows with the literature. Second, the crowders are the readers:
the sign of the spillover from one more paper is derived from the envelope
theorem rather than assumed, and it changes along the payoff curve
(\cref{prop:externality}). Third, the duplication rate is pinned by entry, with
an exact volume elasticity of density built from observables
(\cref{prop:volume}). Fourth, the writing--reading asymmetry, together with a
sufficient statistic for the stepping-on-toes parameter of growth theory
(\cref{prop:sot}). A catch-up cost, the need to read a topic's stock before
contributing \citep{jones2009}, adds a wedge that the city literature lacks:
incumbents' returns rise with density even where newcomers find topics too
crowded (\cref{cor:overcrowd}).

\paragraph{Evidence.} We confront the model with the arXiv record in two ways.
The descriptive evidence (\cref{sec:facts}) comes from 1.49 million papers posted
in 2018--2025, classified into 150 text-based areas. The anchor fact is that volume
opens new ground. From the 95th to the 25th percentile similarity radius, local
density rises with area volume at an elasticity of 0.12--0.19, against 0.53--0.81
if papers kept landing in the same idea space. At the closest-overlap radius the
estimate is 0.035, though it is imprecise (SE 0.118). Returns are hump-shaped in
topic density. Unaware duplicates lose about a fifth of their citations, close to
the 21\% penalty for scooped teams in \citet{hillstein2025}, while aware
follow-ons gain. Once corrected for citation recall, about 60\% of close overlap
is aware.

The quasi-experimental evidence (\cref{sec:causal}) uses arXiv's own rules on a
new panel of 2.2 million arXiv submissions with exact timestamps through
15~October 2026. Four shocks speak to the margins the model relies on.
\begin{itemize}
\item \emph{Entry cost.} In January 2026, arXiv stopped endorsing new submitters
on the strength of an institutional e-mail address alone; Mathematics had moved
in December 2025. The share of papers by first-time submitters fell by
\empEndOther\ percentage points (about \empEndOtherRel) in the affected
sections, net of the year before. The staggered timing matches the rule.
\item \emph{Catch-up cost.} From November 2025, review and position papers in
computer science needed prior peer review. Review-type papers in computer
science halved relative to other archives, more so in categories that had
relied on them. First-time authorship fell more in those categories, though a
smaller pre-trend of the same sign makes this suggestive.
\item \emph{Writing cap.} From October 2026, each person may submit at most two
papers a month. The cap removed about \empCapRemoved\ of papers. Total volume
kept growing at its usual seasonal pace, so per-submitter caps are a weak lever
on volume.
\item \emph{Readership.} A paper's position in the daily listing depends on how
many others submitted earlier that day. We use this to test the
detection channel directly. Position has only a small effect on citations, and
no detectable effect on later overlapping work or on the share of it that cites
the paper, with tight bounds. Overlap is detected by targeted search rather than
by browsing.
\end{itemize}
At the author level, adopters of LLM-style writing raised their output by about
\empLLMallPct\ relative to non-adopters. That is far less than the 36\% of
\citet{kusumegietal2025}, whose design has been criticised for mechanically
generating positive dynamics \citep{renaultetal2026}.

\paragraph{Magnitudes.} The calibrated model is disciplined by the corrected
arXiv moments (\cref{sec:quant}). The reading--writing asymmetry is the
quantitatively robust prediction: a 30\% faster pen raises output per researcher
by \calWriteOut\ log points, while 30\% faster reading raises it by only
\calReadOut. The static welfare cost of crowding is small. Optimal topic choice
would raise output per researcher by \calPlannerGain\ in the baseline calibration
and \calWPlannerGain\ in the within-topic one, and overlap destroys only about
\calErosion\ of credit at mature density. The frictionless benchmark, in which
catch-up is free and topics fill until the whole rent at the peak is dissipated,
would lose \calDissip; it is a counterfactual, not the calibrated economy. The
direction of the inefficiency survives treating all priority credit lost to
overlap as a transfer between racers (\cref{prop:social}). With the measured
volume elasticity, topic crowding explains little of the decline in research
productivity.

\paragraph{Implications.} Two further results follow, and we collect them in
\cref{sec:macro}. Because the stock of literature lags its flow, topics whose
literature grows faster than a threshold rate overshoot their best density.
This survives forward-looking entry: rational entrants overshoot more than
myopic ones once a topic grows fast, because what deters entry is the stock,
and a young topic has little of it (\cref{prop:toofast}). And the
stepping-on-toes parameter equals one minus the product of the volume
elasticity of density and the slope of the entrant's payoff curve.

\paragraph{Related literature.} The paper joins five literatures. The first is
the economics of science and priority
\citep{merton1957,merton1961,dasguptadavid1994,stephan1996,hillstein2025,hillstein2025race,bobtcheffetal2017}
and the philosophy of the division of cognitive labour
\citep{kitcher1990,strevens2003}. We add reading as a choice with two
distinct roles, and derive rather than assume the costs of duplication.
\citet{hillstein2025race} show that priority races lower quality; our cost of
crowding works through detection and catch-up instead, and the two are
complementary. The second is the theory of excess entry, common pools, cities and
clubs
\citep{gordon1954,weitzman1974,mankiwwhinston1986,henderson1974,buchanan1965,hendersonbecker2000,scotchmer2002}
and of patent races \citep{loury1979,dasguptastiglitz1980}. Our resource is
attention to a literature, and the crowders are its readers. The third is the
direction of research
\citep{bryanlemus2017,hopenhaynsquintani2021,carnehlschneider2025,bramoullesaintpaul2010,rzhetskyetal2015,fosteretal2015,myers2020,borjasdoran2012,borjasdoran2015,azoulayetal2019}.
\citet{carnehlschneider2025} model the choice of question by novelty and
feasibility. We model the choice by density, derive the hump, and measure the
margin. \citet{borjasdoran2015} document that researchers move in idea space in
response to crowding, which is the margin our model turns on. The fourth is the
macroeconomics of ideas and of access to knowledge
\citep{romer1990,aghionhowitt1992,jones1995,kortum1997,weitzman1998,bloometal2020,jones2009,ekerdtwu2024,furmanstern2011,murrayetal2016,iariaetal2018,biasimoser2021,evans2008,mccabesnyder2015,agrawaletal2016,parketal2023}.
We give stepping on toes an equilibrium foundation and explain why better access
to the literature need not raise productivity. The fifth is artificial
intelligence in science
\citep{agrawaletal2019,brynjolfssonetal2019,jones2025ai,duedeetal2026,haoetal2026,liangetal2024,kusumegietal2025}.
\citet{jones2025ai} studies AI in a task-based R\&D production function, and
\citet{duedeetal2026} model LLMs as labour-augmenting technology that shifts
effort across projects. We separate writing, entry, reading, search and catch-up
channels by their equilibrium footprints. Methodologically, our text-based
density measures follow \citet{hobergphillips2016} and \citet{kellyetal2021}, and
the listing-position design builds on \citet{haqueginsparg2009,haqueginsparg2010}
and \citet{feenbergetal2017}.

\Cref{sec:model} sets up the model and \cref{sec:payoff} derives the payoff
curve. \Cref{sec:frictionless,sec:catchup} characterise equilibrium, welfare and
social value. \Cref{sec:ai} states the writing--reading asymmetry.
\Cref{sec:facts,sec:causal} present the descriptive and quasi-experimental
evidence, \cref{sec:quant} quantifies the model, \cref{sec:macro} collects
the implications for fast-growing topics and growth, \cref{sec:predictions}
lists out-of-sample predictions, and \cref{sec:conclusion} concludes.
